\section{Математическая постановка задачи}

Рассмотрим геораспределённую систему межкластерной репликации данных,
функционирующую в условиях сетевых задержек, отказов узлов и каналов связи,
а также возможных сетевых и инфраструктурных атак. Требуется формально
описать модель системы и сформулировать задачу выбора параметров репликации,
направленную на обеспечение целостности, доступности и согласованности данных.

\subsection{Основные множества и топология системы}

Пусть
$\mathcal{C} = \{c_1, c_2,\dots, c_K\}$
--- множество кластеров геораспределённой системы, где $K$ --- общее число кластеров.

В дальнейшем рассматривается взаимодействие пары кластеров
$\{c_a,c_p\}\subseteq\mathcal{C}$ в архитектуре ``активный--пассивный'';
прочие кластеры (при наличии) не участвуют в рассматриваемом канале репликации.

Взаимодействие между кластерами задаётся ориентированным графом~\ref{eq:G}:
\begin{equation}
\mathcal{G} = (\mathcal{C}, \mathcal{E}),
\label{eq:G}
\end{equation}
где $\mathcal{E} \subseteq \mathcal{C} \times \mathcal{C}$ ---
множество ориентированных рёбер (каналов связи) между кластерами.

Элемент
$(c_i, c_j) \in \mathcal{E}$
означает наличие канала передачи данных из кластера $c_i$ в кластер $c_j$.
Отсутствие пары $(c_i,c_j)$ в множестве $\mathcal{E}$ означает невозможность
непосредственного обмена сообщениями между указанными кластерами.
Таким образом, граф $\mathcal{G}$ формально описывает сетевую топологию
геораспределённой системы.

Для каждого $c_i \in \mathcal{C}$ определено множество узлов кластера
$
\mathcal{N}_i = \{1,2,\dots,n_i\},
$
где $n_i$ --- число узлов в кластере $c_i$.

\subsection{Модель данных и операций}

Будем считать, что система функционирует во времени
$t \in [0,+\infty)$,
где $t$ --- непрерывная временная переменная.

Пусть
$\mathcal{X} = \{x_1, x_2, \dots, x_M\}$
--- множество реплицируемых объектов, где $M$ --- их количество.

В каждом кластере $c_i\in\mathcal{C}$ выполняется последовательность операций
$\mathcal{O}_i = \{o_{i,1}, o_{i,2}, \dots\}$,
где индекс $i$ указывает на принадлежность операции кластеру $c_i$,
а индекс $k$ --- на порядковый номер операции в данном кластере
в порядке её возникновения во времени.
Каждая операция задаётся кортежем~\ref{eq:op}:
\begin{equation}
o_{i,k} = (\operatorname{op}, x, v, \tau_{i,k}),
\label{eq:op}
\end{equation}
где
\begin{itemize}
\item $\operatorname{op} \in \{\mathrm{read}, \mathrm{write}\}$ --- тип операции,
\item $x \in \mathcal{X}$ --- объект доступа,
\item $v$ --- значение (для операции записи),
\item $\tau_{i,k} \in [0,+\infty)$ --- момент времени выполнения операции.
\end{itemize}

Заметим, что значения $\tau_{i,k}$ используются как абстрактные временные метки
внутри модели и задают порядок операций в одном кластере.

Выделим подпоследовательность операций записи:
\[
\mathcal{O}_i^{\mathrm{w}}=\{o_{i,k}\in\mathcal{O}_i \mid \operatorname{op}=\mathrm{write}\}.
\]

Операции записи порождают журнал репликации кластера.
Журнал представим в виде последовательности записей
\begin{equation}
\mathcal{L}_i = (e_{i,1}, e_{i,2}, \dots),
\label{eq:Li}
\end{equation}
где $e_{i,\ell}$ --- $\ell$-я запись журнала, а индекс $\ell \in \mathbb{N}$
является её логическим номером (порядковым номером в журнале).

Связь между операциями записи и записями журнала зададим отображением
\begin{equation}
\varphi_i:\mathcal{O}_i^{\mathrm{w}} \to \mathcal{L}_i,
\label{eq:phi_i}
\end{equation}
которое каждой операции записи сопоставляет одну запись журнала.
Тем самым журнал $\mathcal{L}_i$ является упорядоченным результатом применения
$\varphi_i$ ко всем операциям записи в порядке их возникновения в кластере.

Для каждого момента времени $t$ обозначим через $\mathcal{L}_i(t)$ префикс журнала,
сформированный к моменту времени $t$:
\[
\mathcal{L}_i(t) = (e_{i,1}, e_{i,2}, \dots, e_{i,a_i(t)}),
\]
где $a_i(t)\in\mathbb{N}$ --- число записей, сформированных (и зафиксированных)
в кластере $c_i$ к моменту времени $t$.

Межкластерная репликация по каналу $(c_i,c_j)\in\mathcal{E}$
описывается оператором доставленного префикса:
\begin{equation}
\mathcal{R}_{i\to j}(t): \mathcal{L}_i(t) \to \mathcal{L}_j(t),
\label{eq:Rij_time}
\end{equation}
где $\mathcal{R}_{i\to j}(t)$ формально означает ту часть записей из $\mathcal{L}_i(t)$,
которая была доставлена и принята кластером $c_j$ к моменту времени $t$.

\subsection{Модель сети, задержек и отказов}

Функционирование рассматриваемой системы
зависит от случайных сетевых и инфраструктурных факторов,
влияющих на доставку репликационных сообщений
и работоспособность узлов.
Будем считать, что существует вероятностное пространство
\[
(\Omega,\mathcal{F},\mathbb{P}),
\]
где:

\begin{itemize}
\item $\Omega$ — множество возможных сценариев функционирования системы,
включающих различные комбинации сетевых задержек,
потерь сообщений и отказов узлов;
\item $\mathcal{F}$ — совокупность событий,
связанных с этими сценариями;
\item $\mathbb{P}$ — вероятность реализации соответствующего сценария.
\end{itemize}

Таким образом, случайность в модели связана
не с логикой работы протокола,
а с внешними воздействиями:
сетевыми задержками,
потерями пакетов и отказами оборудования.

Для каждого канала $(c_i,c_j)\in\mathcal{E}$ введём следующие случайные процессы:

\begin{enumerate}

\item
$\Delta_{ij}(t) \ge 0$
--- процесс задержки доставки сообщения
из кластера $c_i$ в кластер $c_j$ в момент времени $t$.

\item
$D_{ij}(t) \in \{0,1\}$
--- индикатор доставки сообщения,
где $D_{ij}(t)=1$ соответствует успешной доставке репликационного пакета
из кластера $c_i$ в кластер $c_j$, а $D_{ij}(t)=0$ --- потере пакета.

\item
$S_{i,u}(t) \in \{\mathrm{up}, \mathrm{down}\}$
--- процесс состояния узла $u \in \mathcal{N}_i$,
где $\mathrm{up}$ означает корректную работу узла,
а $\mathrm{down}$ --- его отказ.

\end{enumerate}

\subsection{Модель атак}

Сетевые атаки моделируются как воздействие на каналы связи
между кластерами системы.
Пусть $(c_i, c_j) \in \mathcal{E}$.
Противник может оказывать влияние на случайные процессы
$\Delta_{ij}(t)$ и $D_{ij}(t)$, определённые в модели сети.

Допускаются следующие типы воздействий:

\begin{enumerate}
\item \textbf{Увеличение задержки доставки сообщений.}
Противник может искусственно увеличивать время передачи сообщения
между кластерами. Формально, это описывается преобразованием~\ref{eq:delay_attack}:
\begin{equation}
\Delta_{ij}(t) \rightarrow \Delta_{ij}(t) + \delta_{ij}(t),
\label{eq:delay_attack}
\end{equation}
где $\delta_{ij}(t) \ge 0$ --- дополнительная задержка,
вносимая атакой.
Это означает, что фактическое время доставки сообщения
становится больше номинального значения задержки канала.

\item \textbf{Принудительная потеря сообщений.}
Противник может инициировать потерю передаваемых сообщений.
Формально это описывается равенством
$
D_{ij}(t) = 0,
$
что означает отсутствие доставки репликационного пакета
в момент времени $t$.

\end{enumerate}

Инфраструктурные атаки моделируются как отказ части узлов кластеров.

Пусть
$
\mathcal{F}_i(t) \subseteq \mathcal{N}_i,
$
где $\mathcal{F}_i(t)$ --- множество узлов кластера $c_i$,
находящихся в состоянии отказа в момент времени $t$.
Предполагается ограничение на мощность воздействия:
$
|\mathcal{F}_i(t)| \le f_i,
$
где $f_i$ --- максимально допустимое число одновременно
отказавших узлов в кластере $c_i$. Тогда состояние узла $u \in \mathcal{N}_i$ определяется как
\[
S_{i,u}(t) =
\begin{cases}
\mathrm{up}, & u \notin \mathcal{F}_i(t),\\
\mathrm{down}, & u \in \mathcal{F}_i(t).
\end{cases}
\]

Таким образом, инфраструктурные атаки приводят к снижению
работоспособности кластеров и могут влиять на доступность
и корректность межкластерной репликации.

В соответствии с формулировкой задачи необходимо обеспечить
целостность, доступность и согласованность данных
в условиях сетевых задержек и отказов узлов.

\subsection{Критерий целостности данных}

Перед формулировкой критериев введём вектор управляемых параметров системы
\[
\theta=
\bigl(
m,\;
\tau_{\mathrm{fo}},\;
\tau_{\mathrm{el}},\;
\tau_{\mathrm{sp}},\;
r
\bigr),
\]
где $m$ — минимальное число исправных узлов активного кластера,
$\tau_{\mathrm{fo}}$ — время ожидания смены активного кластера,
$\tau_{\mathrm{el}}$ — время ожидания внутреннего лидерства,
$\tau_{\mathrm{sp}}$ — время ожидания подтверждения состояния,
$r$ — число повторных попыток передачи репликационного сообщения.

Целостность данных при межкластерной репликации будем понимать как
сохранение порядка и содержимого записей журнала при доставке в пассивный кластер,
а также отсутствие вставок и пропусков в доставленной части.

Поскольку журналы $\mathcal{L}_i(t)$ и $\mathcal{L}_j(t)$ являются последовательностями,
введём отношение ``быть префиксом''.
Для двух последовательностей $A$ и $B$ будем писать $A \preceq B$,
если $A$ является префиксом $B$ (то есть $B$ начинается с элементов $A$
в том же порядке).

Для фиксированной конфигурации параметров $\theta$
обозначим через $\mathcal{R}_{i\to j}(t;\theta)$
доставленную к моменту времени $t$ часть журнала источника.
Данная зависимость от $\theta$ возникает потому,
что параметры таймаутов и число повторных попыток передачи
влияют на процесс доставки и применения записей.

Требование целостности формулируется как выполнение двух условий:

\begin{enumerate}
\item \textbf{Префиксность доставки:}
доставленные записи образуют префикс журнала источника,
то есть
\[
\mathcal{R}_{i\to j}(t;\theta) \preceq \mathcal{L}_i(t);
\]

\item \textbf{Префиксность применения:}
журнал получателя содержит доставленные записи как префикс,
то есть
\[
\mathcal{R}_{i\to j}(t;\theta) \preceq \mathcal{L}_j(t).
\]
\end{enumerate}

Нарушением целостности будем называть событие

\[
\mathsf{Int}_{i\to j}(t;\theta)
=
\neg\Bigl(
\mathcal{R}_{i\to j}(t;\theta) \preceq \mathcal{L}_i(t)
\;\wedge\;
\mathcal{R}_{i\to j}(t;\theta) \preceq \mathcal{L}_j(t)
\Bigr).
\]

Поскольку система функционирует во времени,
оценку целостности удобно задавать в усреднённой по времени форме.
Требование целостности формулируется как

\begin{equation}
\frac{1}{T}
\int_{0}^{T}
\mathbb{P}\!\left(\mathsf{Int}_{i\to j}(t;\theta)\right)dt
\le
\varepsilon_I,
\qquad
\forall (c_i,c_j)\in\mathcal{E},
\label{eq:integrity_final}
\end{equation}

где $\varepsilon_I>0$ — допустимый уровень риска нарушения целостности.

\subsection{Критерий доступности данных}

В системе предполагается наличие активного кластера,
обслуживающего операции чтения и записи,
и пассивного кластера,
обслуживающего только операции чтения.

Для фиксированной конфигурации параметров $\theta$
обозначим через

\[
A_i^{\mathrm{r}}(t;\theta) \in \{0,1\}
\]

случайную величину, характеризующую возможность выполнения операции чтения
в кластере $c_i$ в момент времени $t$.

Аналогично, для активного кластера введём

\[
A_a^{\mathrm{w}}(t;\theta) \in \{0,1\},
\]

характеризующую возможность выполнения операции записи.

Поскольку состояние системы определяется
случайными задержками $\Delta_{ij}(t)$,
случайными отказами узлов,
а также выбранной конфигурацией параметров $\theta$,
величины
$A_i^{\mathrm{r}}(t;\theta)$
и
$A_a^{\mathrm{w}}(t;\theta)$
являются случайными.

Для количественной оценки доступности используется математическое ожидание

\[
M\big(A_i^{\mathrm{r}}(t;\theta)\big),
\qquad
M\big(A_a^{\mathrm{w}}(t;\theta)\big),
\]

которое интерпретируется как вероятность успешного обслуживания
соответствующего запроса в момент времени $t$.

Поскольку состояние системы меняется во времени,
необходимо учитывать динамический характер отказов.
Для этого рассматривается средняя доля времени
успешного функционирования на интервале $[0,T]$.

Требования доступности формулируются следующим образом:

\begin{equation}
\begin{cases}
\displaystyle
\frac{1}{T}
\int_{0}^{T}
M\big(A_a^{\mathrm{w}}(t;\theta)\big)\, dt
\ge
1 - \varepsilon_A^{\mathrm{w}}, \\[1.2em]

\displaystyle
\frac{1}{T}
\int_{0}^{T}
M\big(A_i^{\mathrm{r}}(t;\theta)\big)\, dt
\ge
1 - \varepsilon_A^{\mathrm{r}},
\qquad
i \in \{a,p\}.
\end{cases}
\label{eq:availability_system}
\end{equation}

где
$\varepsilon_A^{\mathrm{w}}$ и
$\varepsilon_A^{\mathrm{r}}$
— допустимые уровни недоступности
для операций записи и чтения соответственно.

\subsection{Критерий согласованности данных}

В архитектуре ``активный--пассивный''
согласованность естественно измерять лагом репликации,
то есть отставанием пассивного кластера по журналу записей.

Пусть $\mathcal{L}_a(t)$ — префикс журнала активного кластера,
сформированный к моменту времени $t$,
а $\mathcal{L}_p(t)$ — префикс журнала пассивного кластера.

Обозначим через $|\mathcal{L}_i(t)|$
длину соответствующего префикса журнала.

Для фиксированной конфигурации параметров $\theta$
лаг репликации определим как

\begin{equation}
L(t;\theta)
=
|\mathcal{L}_a(t)|
-
|\mathcal{L}_p(t)|.
\label{eq:log_lag}
\end{equation}

Величина $L(t;\theta)$ зависит от параметров конфигурации $\theta$,
поскольку параметры таймаутов и число повторных передач
влияют на скорость доставки и применения записей.

По определению $L(t;\theta)\ge0$,
так как пассивный кластер получает записи из активного.

Для количественной оценки согласованности
используется математическое ожидание лага
$M(L(t;\theta))$.

Требование согласованности формулируется в виде ограничения
на средний лаг репликации:

\begin{equation}
\frac{1}{T}
\int_{0}^{T}
M\bigl(L(t;\theta)\bigr)\,dt
\le
\Delta_L,
\label{eq:consistency_log}
\end{equation}

где $\Delta_L$ — допустимый уровень среднего лага репликации.

\subsection{Итоговая формулировка задачи}

Параметр $m$ определяет минимальное число исправных узлов,
при котором активный кластер считается работоспособным.

Параметры
$\tau_{\mathrm{fo}}, \tau_{\mathrm{el}}, \tau_{\mathrm{sp}}$
характеризуют временные ожидания различных процедур
обнаружения отказов и восстановления состояния системы.
Параметр $r$ определяет число повторных попыток передачи
репликационного сообщения.

Обозначим через $\Theta$
множество всех допустимых конфигураций параметров $\theta$.
Для оценки влияния параметра $m$ предположим,
что каждый узел активного кластера находится в состоянии отказа
с вероятностью $p$ независимо от других узлов.
Тогда число исправных узлов имеет биномиальное распределение,
и вероятность сохранения активности кластера равна

\[
P_{\mathrm{act}}(m)
=
\sum_{k=m}^{n_a}
\binom{n_a}{k}
(1-p)^k p^{\,n_a-k}.
\]

Аналогично,
если вероятность успешной доставки сообщения за одну попытку равна $p_d$,
то вероятность успешной доставки при $r$ попытках равна

\[
P_{\mathrm{succ}}(r)
=
1-(1-p_d)^r.
\]

Для сравнения различных конфигураций параметров
введём функцию эксплуатационных издержек
$
J:\Theta\to\mathbb{R}_{+}.
$
Эксплуатационные характеристики системы
(целостность, доступность и согласованность)
задаются в виде ограничений на допустимые конфигурации параметров.
Функция эксплуатационных издержек используется
для сравнения различных допустимых конфигураций
и выбора наименее затратной из них.

Функционал определим следующим образом:

\begin{equation}
J(\theta)
=
\bigl(1-P_{\mathrm{act}}(m)\bigr)
+
\frac{\tau_{\mathrm{el}}}{\tau_{\mathrm{fo}}}
+
\frac{\tau_{\mathrm{sp}}}{\tau_{\mathrm{fo}}}
+
\bigl(1-P_{\mathrm{succ}}(r)\bigr),
\label{eq:J_final}
\end{equation}

где
$
\theta=
\bigl(
m,\;
\tau_{\mathrm{fo}},\;
\tau_{\mathrm{el}},\;
\tau_{\mathrm{sp}},\;
r
\bigr),
$

Первое слагаемое отражает риск потери активности кластера,
второе и третье характеризуют относительную стоимость временных ожиданий,
а четвёртое — риск неуспешной доставки репликационного сообщения.

Все слагаемые являются безразмерными величинами,
что позволяет корректно агрегировать их
в едином функционале издержек.

Определим множество допустимых конфигураций параметров

\[
\Theta_{\mathrm{adm}}
=
\left\{
\theta\in\Theta
\;\middle|\;
\begin{aligned}
&\frac{1}{T}\int_0^T
\mathbb{P}\!\left(\mathsf{Int}_{i\to j}(t;\theta)\right)dt
\le\varepsilon_I,
\qquad \forall (c_i,c_j)\in\mathcal{E},\\[0.4em]
&\frac{1}{T}\int_0^T
M\big(A_a^{\mathrm{w}}(t;\theta)\big)dt
\ge 1-\varepsilon_A^{\mathrm{w}},\\[0.4em]
&\frac{1}{T}\int_0^T
M\big(A_i^{\mathrm{r}}(t;\theta)\big)dt
\ge 1-\varepsilon_A^{\mathrm{r}},
\qquad i\in\{a,p\},\\[0.4em]
&\frac{1}{T}\int_0^T
M\big(L(t;\theta)\big)dt
\le \Delta_L
\end{aligned}
\right\},
\]

где 
$
\theta=
\bigl(
m,\;
\tau_{\mathrm{fo}},\;
\tau_{\mathrm{el}},\;
\tau_{\mathrm{sp}},\;
r
\bigr),
$

Требуется определить конфигурацию параметров

\[
\theta^{*}
=
\arg\min_{\theta\in\Theta_{\mathrm{adm}}}
J(\theta),
\]

то есть найти такие значения
$m,\tau_{\mathrm{fo}},\tau_{\mathrm{el}},\tau_{\mathrm{sp}},r$,
которые обеспечивают выполнение требований
целостности, доступности и согласованности
и одновременно минимизируют эксплуатационные издержки системы.
